% Introduction. Related work (former Section 2) is folded in here, in the JFM
% idiom (no standalone "Related work" section). The SIGReg-vs-PLDM "second paper"
% that headed Section 1.2 of the original draft has been removed; the regulariser
% choice is now a methods detail (Sec. 3) and an appendix. Reframed 2026-06-08 for
% the fully UNCONDITIONED model: no gust parameters enter the encoder or the
% predictor. Contributions are reframed around the representational wake closure
% the predictive objective produces, with the rollout shown but not headlined as
% forward closure, and the world-model reading kept as motivation only.

A discrete vortex gust striking a post-stall airfoil does not merely perturb the
force coefficients: it reorganises the separated wake. The lift excursion it
produces is the surface signature of leading-edge-vortex roll-up, growth, and
shedding, the dominant unsteady mechanism in these encounters
\citep{eldredge2019leading}, together with the trailing-edge and wake structures
it leaves behind, and it is set by the timing of the vortex impact relative to the
airfoil's own shedding, the wall-normal offset of the vortex core, and the gust
strength and core size. Such extreme gusts arise wherever the gust velocity is
comparable to or larger than the flight speed, for small uncrewed and micro air
vehicles in urban canyons, mountainous terrain, and the wakes of buildings and
ships \citep{jones2022,mohamed2023gusts}: the gust ratio $G$, the ratio of the
characteristic gust velocity to the free-stream velocity, then exceeds unity, and
conventional gust models do not capture the resulting large, fast load transients
\citep{fukami2024control,fukami2025prf}. Even at the moderate chord Reynolds number
$\vRe = 5000$, the post-stall flow at $\alpha = 14^\circ$ is massively
separated, and the gust-induced disturbance interacts nonlinearly with the
natural shedding cycle \citep{fukami2025prf}; the vortical organisation of the
encounter is moreover similar across laminar and turbulent regimes
\citep{odaka2026}, so this case is representative of the wider extreme-aerodynamic
family. Resolving the full parametric
envelope by simulation is expensive, so reduced-order models (ROMs) of the
encounter-to-encounter and case-to-case variability are an active line of work.
A reduced model for such an encounter must do more than reproduce a vorticity
snapshot or the instantaneous force: it must retain the wake variables that
determine the post-impact transient and organise them into a state on which a
forward model can be queried without leaving the data manifold. This is the
predictive-state problem that organises the present study.

Reduced-order models of such flows form a ladder of increasing nonlinearity.
Energy-optimal linear bases, proper orthogonal decomposition
\citep{berkooz1993pod} and dynamic mode decomposition \citep{schmid2010dmd}, are
compact and interpretable but limited for strongly nonlinear separation; their
Koopman-operator reading \citep{lusch2018deep} nonetheless makes a linear latent a
principled comparator rather than a strawman. Nonlinear autoencoders compress
further at matched latent dimension \citep{murata2020nonlinear,lee2020model}, and
coupling an
autoencoder to a learned latent dynamics forecasts better than dynamics confined
to a linear subspace \citep{linot2023dynamics}. Read as a design ladder, these
methods progress from field reconstruction and energy compression, through
observable readability, toward forward-state usability, which is the property a
reduced state needs to serve a predictor; the extreme-aerodynamic reduced models we
build on sit at the nonlinear, observable-augmented rung, and the question we pursue
is what carries a state the rest of the way.

Recent ROMs for separated and extreme-aerodynamic flows, increasingly built with
machine learning \citep{brunton2020machine}, share a common
ingredient: a reconstruction-trained autoencoder whose bottleneck is shaped by
an auxiliary aerodynamic observable. Fukami \& Taira \citeyearpar{fukami2023} trained
a convolutional autoencoder jointly with a lift-prediction head and found that
the auxiliary loss reduced the intrinsic dimension of the discovered manifold
for extreme vortex-gust flows from roughly ten to three. The same
observable-augmentation idea was extended across multi-source turbulent data
\citep{fukami2025}, and was used by Solera-Rico et al. \citeyearpar{solera2024} with a
$\beta$-variational autoencoder and a transformer ROM on a vortex-gust airfoil
dataset of the same family considered here. Across this lineage the
reconstruction term anchors the latent geometry to the instantaneous flow
field, and an observable-prediction term orients it along aerodynamically
relevant directions. A recurring theme of the geometric-analysis literature,
however, is that the latent coordinates of an autoencoder are unique only up to
a diffeomorphism, so reconstruction leaves the geometry of the latent largely
unconstrained, which complicates any downstream task that relies on that
geometry, including dynamics, interpolation, and comparison
\citep{tran2026,kelshaw2024}. Observable augmentation solves only one part of the
predictive-state problem: it tells the encoder which aerodynamic quantity to
represent, but it does not by itself fix whether the latent coordinates are
suitable for time advancement. A reduced state that is linearly readable at a single
instant may still be unusable once an autoregressive predictor queries the points
its own rollout produces.

A parallel and fast-growing line reads these flows through the information their
reduced states carry rather than their energy. Transfer-entropy causality among
proper-orthogonal-decomposition modes ranks which modes drive the formation of
coherent structures \citep{martinezsanchez2023}, and an informative/non-informative
decomposition splits a field into the part that carries information about a
target's future and a residual that carries none \citep{arranz2024}. This
informative-decomposition idea has very recently been brought to extreme
vortex-gust airfoil interactions of the present family, extracting the field
structures informative about the future lift \citep{koshikawa2026,fukamiaraki2026},
alongside time-dependent modal bases for the same interaction
\citep{zhong2025,zamaniashtiani2025} and latent-causality attributions in physical
space \citep{cremades2026xcal}. Our study is complementary: rather than mapping the
informative or causal modes of a single representation, we ask which encoder
objective produces a latent whose forward evolution closes the full set of
observables, the wake included, and we use the information lens only as a held-out
diagnostic of the resulting latent (\S\ref{sec:disc_mechanism}).

A complementary route to a reduced state is to drop the reconstruction term
altogether and to train the encoder and a predictor end to end on
latent-space forecasting, with an explicit anti-collapse regulariser in place
of the decoder. This is the joint-embedding predictive architecture (JEPA)
\citep{lecun2022}: an encoder maps the input to a representation, a predictor
advances that representation in time, and the loss never reconstructs the
field. The encoder is thereby trained in the same coordinates in which the
predictor will later operate, which makes JEPA a useful instrument for testing
whether predictive training produces a latent state that remains usable when rolled
forward, rather than merely readable when encoded from the DNS field; the predictor
also cannot exploit pixel-level shortcuts, because no field-space signal reaches it. The image and video
instantiations \citep{ijepa,vjepa2} and the recent from-pixels variants that
remove the moving-target encoder in favour of a characteristic-function
regulariser \citep{lewm,lejepa} or a variance-covariance objective
\citep{pldm,vicreg} have so far been evaluated on gridworld and toy visual
domains. Their behaviour on a low-intrinsic-dimension physics dataset, where
the data sit on a roughly five- to ten-dimensional manifold, is uncharacterised.

Reconstruction-trained nonlinear encoders compress aggressively but yield latent
dynamics that are harder to forecast than an energy-optimal linear basis, a
compactness-versus-forecast tradeoff documented for controlled wake flows in a
companion study under review \citep{solerarico_compactness_underreview}. We show
this tradeoff is a consequence of the reconstruction objective: a predictive
latent recovers nonlinear compactness without surrendering forecast stability.
The division of labour with that companion study
\citep{solerarico_compactness_underreview} is explicit: there, controlled
canonical wakes and the compactness-versus-forecast tradeoff; here, the
parametric vortex gust, the latent-drift mechanism we identify as the common
cure, and its geometric and topological characterisation.

The model we adopt takes no gust parameters: $(G, D, Y)$ enter neither the encoder
nor the predictor. The encoder is a pure state map and the predictor advances the
latent from its own history alone. The world-model framework of
\citet{lecun2022,ha2018world} motivates this split as an analogy and not a property
we claim, since the latent dynamics function is the object a controller would plan
against, and a recent gust-airfoil instance couples a physics-augmented autoencoder
to a latent dynamics model \citep{liu2025modelbased}. We treat that reading as
motivation: the rollout is shown as a consistent latent trajectory rather than
advanced as a validated forecast, and its interventional or causal interpretation is
left for future work (\S\ref{sec:discussion}). What this leaves is the
property our study turns on: an established reconstruction-trained autoencoder and an
energy-optimal proper orthogonal decomposition basis, lacking wake supervision, do
not keep the dynamically relevant wake observables linearly readable, and even a
wake-readable state is not automatically usable once a forward model advances it.
The drift mechanism of \S\ref{sec:res_drift} and the division of labour we identify
are the empirical consequences.

The question this paper addresses is therefore not which encoder reconstructs
the flow best, nor even which produces a reduced state from which the wake is
readable at a single instant, but which produces a state that is also usable once a
forward model advances it. We treat this as a fluid-mechanics question with a
geometric answer, and find a division of labour: wake-observable supervision makes
the wake readable, an anti-collapse regulariser keeps the latent in distribution
under rollout, and the predictive objective makes the supervised state dynamically
trackable. The joint-embedding predictive architecture is the instrument we use to
test this principle, not the subject of the paper. We make three contributions.

First, we formulate the vortex-gust encounter as a predictive-state benchmark. The
primary endpoint is the wake enstrophy at a fixed post-impact readout ($H = 16$),
because it summarises the leading-edge-vortex roll-up and wake reorganisation that
the scalar force can miss; all encoders, a predictive one (JEPA), a reconstructive
one in the Fukami lineage, and a POD basis, are compared at matched latent dimension
under a single shared probe and autoregressive predictor (\S\ref{sec:methods}). The
question is then not field-reconstruction quality but whether the reduced state
carries the wake variables a forward model needs (\S\ref{sec:res_closure}).

Second, we separate the ingredients of a usable predictive state
(\S\ref{sec:res_controls}). Wake-observable supervision supplies the instantaneous
readability: on held-out fields the wake-supervised states make the wake enstrophy
linearly readable ($R^2 = \NumReprWakeJepaTf$) where the no-wake reconstructive and
linear baselines are at or below zero, but an objective-free supervised encoder
reads it as well, and a matched reconstructive control reads it comparably once
retrained out of sample, so the readability is the supervision's, not the
objective's. Anti-collapse regularisation supplies an in-distribution geometry: the
supervision-only and predictive rollouts both stay on the encoded manifold, where
the reconstructive rollout departs along directions of negligible encoded variance
(\S\ref{sec:res_drift}). The predictive objective supplies the dynamic usability: its
rollout stays closer to the encoded simulation trajectory and forecasts the wake
better than the objective-free and reconstructive alternatives.

Third, we interpret the resulting state physically and observationally
(\S\ref{sec:res_flow_physics}, \S\ref{sec:res_observability}). The predictive state
keeps the undisturbed shedding orbit as a single loop under a metric that removes
coordinate scale, carries the wake response across the tested readout horizons, and
is the most recoverable of the three from sparse pre-impact wall pressure. These
results motivate the state as an input to a future estimator or controller; the
present paper does not claim a validated forecast or a closed-loop controller.
